\documentclass[10pt]{article}
\usepackage[a4paper,margin=0.7in]{geometry}
\usepackage{amsmath,amssymb,amsfonts,mathtools}
\setlength{\jot}{4pt}
\usepackage{xcolor}
\usepackage{hyperref}
\setlength{\parindent}{0pt}
\setlength{\parskip}{5pt}
\newcommand{\Lc}{\mathcal L}
\newcommand{\Jc}{\mathcal J}
\newcommand{\Bc}{\mathcal B}
\newcommand{\Ac}{\mathcal A}
\newcommand{\Nn}{\mathcal N}
\newcommand{\Ph}{\hat P}
\title{Final-state $g\to gg$ contribution (Group II, Row IV):\\
the complete analysis in coordinate space}
\author{}
\date{}
\begin{document}
\maketitle
\vspace{-2.5em}

\begin{abstract}
We redo, step by step and entirely in coordinate space, the analysis of the note
\emph{``Divergence structure of the final-state $g\to gg$ contribution''}: exact reduction, the
target-rapidity (JIMWLK) and projectile-rapidity (BFKL) endpoints, the collinear (DGLAP) pole and
its plus prescription, the $\beta_0\,\delta(1-\xi)$ completion and the coordinate-space running
coupling, the net-colour large-distance IR divergence, and the finite remainder. The only
momentum-space input is the closed form of one coordinate kernel $H^{mm'}(s)$, the
coordinate image of the $\overline{\rm MS}$-regulated splitting kernel. All results agree with
the note, and the endpoint coefficients agree with the ones obtained earlier in coordinate
space.
\end{abstract}

\section{Starting point, notation and scheme}
The contribution is
\begin{align}\label{eq:start}
\frac{d^{3}N}{d^{3}k}= &\frac{1}{(2\pi)^{3}}\frac{g^{4}N_{c}}{4\pi^{4}}\frac{1}{k^{+}}\int_{\xi_V}^{\xi_\Lambda}\frac{d\xi}{2\pi}\int_{z,w,w',x'}e^{-ik\cdot(w'-w)}\int_{x,y,y'}\frac{(w-z)^{m}(x-y)^{k}(w'-z)^{m'}(x'-y')^{k'}}{(w-z)^{2}(x-y)^{2}(w'-z)^{2}(x'-y')^{2}}\notag\\
&\times\delta^{(2)}\big[y-z-\xi(w-z)\big]\,\delta^{(2)}\big[y'-z-\xi(w'-z)\big]\,T_{mkm'k'}(\xi)\,
\left[U^{ce'}(y')-U^{ce'}(x')\right]\left[U^{ce}(y)-U^{ce}(x)\right]\rho^{e'}(x')\rho^{e}(x),
\end{align}
with $T_{mkm'k'}=d_\perp\xi(1-\xi)\delta_{km}\delta_{k'm'}+2(\delta_{mk'}\delta_{km'}-\delta_{km}\delta_{k'm'})+c(\xi)\delta_{mm'}\delta_{kk'}$,
$c(\xi)=\frac{\xi}{1-\xi}+\frac{1-\xi}{\xi}$, $\xi_V=\frac{k^{+}}{\vee}$,
$\xi_\Lambda=\frac{k^{+}}{k^{+}+\Lambda}$.

\textbf{Notation (coordinate space).} The LO amplitude and its bilinears:
\begin{align}
A^{c}_{i}(y)&=\int_x\frac{(x-y)_i}{(x-y)^{2}}\big[U(y)-U(x)\big]^{ce}\rho^{e}(x),\qquad
\Ac(y,y')\equiv A^{c}(y)\!\cdot\!A^{c}(y'),\qquad
\Ac_{\hat s}(y,y')\equiv\big(\hat s\!\cdot\!A^{c}(y)\big)\big(\hat s\!\cdot\!A^{c}(y')\big),
\end{align}
with $\hat s=(y'-y)/|y'-y|$ and the ordering $\rho^{e'}(x')\rho^{e}(x)$ of Eq.~\eqref{eq:start}
kept throughout. Explicitly
$\Ac(y,y')=\int_{x,x'}\frac{(x-y)\cdot(x'-y')}{(x-y)^{2}(x'-y')^{2}}[U^{ce'}(y')-U^{ce'}(x')][U^{ce}(y)-U^{ce}(x)]\rho^{e'}(x')\rho^{e}(x)$.
The LO spectrum at transverse momentum $p$ and the prefactor are
\begin{align}
\Lc(p)\equiv\int_{y,y'}e^{-ip\cdot(y'-y)}\,\Ac(y,y'),\qquad
\Nn\equiv\frac{1}{(2\pi)^{3}}\frac{g^{4}N_{c}}{8\pi^{5}}\frac{1}{k^{+}}=\frac{\bar C}{k^{+}},\qquad
\Ph(\xi)=\frac{\xi}{1-\xi}+\frac{1-\xi}{\xi}+\xi(1-\xi).
\end{align}

\textbf{Scheme (as in the note).} $d=2-2\epsilon$, $d_\perp=d$; the collinear divergence is an
IR divergence of the \emph{relative} transverse momentum of the pair, so $\epsilon<0$;
$\overline{\rm MS}$ measure $\mu^{2\epsilon}(e^{\gamma_E}/4\pi)^{\epsilon}(2\pi)^{2\epsilon}\int d^{d}\ell$;
scaleless virtual integrals are set to zero ($1/\epsilon_{\rm UV}=1/\epsilon_{\rm IR}$). In
coordinate space this means: the relative coordinate $r$ of the two daughters (conjugate to
$\ell$) is the dimensionally continued variable. Rapidity divergences use the cutoffs
$\xi_V,\xi_\Lambda$.

\section{Exact reduction in coordinate space}
\textbf{Change of variables.} Use the $\delta$-functions to trade $(z,w,w')$ for the parents
$(y,y')$ and the daughters' relative coordinate $r\equiv w-z$:
\begin{align}
z=y-\xi r,\qquad w=y+(1-\xi)r,\qquad w'-z=r+s,\qquad s\equiv\frac{y'-y}{\xi}=w'-w.
\end{align}
The Jacobian is $|\det\partial(z,w,w')/\partial(y,y',r)|=1/\xi$ per transverse component, i.e.\
$1/\xi^{2}$ in two dimensions. The phase becomes $e^{-ik\cdot(w'-w)}=e^{-ik\cdot(y'-y)/\xi}$, the
two emission kernels become $\frac{r^{m}}{r^{2}}$ and $\frac{(r+s)^{m'}}{(r+s)^{2}}$, and
$A(y)A(y')$ does not depend on $r$. Hence, exactly,
\begin{align}\label{eq:Ixi}
\frac{d^{3}N}{d^{3}k}=\frac{\bar C}{k^{+}}\int_{\xi_V}^{\xi_\Lambda}\!d\xi\,I(\xi),\qquad
I(\xi)=\frac{1}{\xi^{2}}\int_{y,y'}e^{-ik\cdot(y'-y)/\xi}\;T_{mkm'k'}(\xi)\,H^{mm'}\Big(\frac{y'-y}{\xi}\Big)\,A^{c}_{k}(y)A^{c}_{k'}(y'),
\end{align}
with the \emph{daughter kernel}
\begin{align}
H^{mm'}(s)\equiv\int_{r}\frac{r^{m}}{r^{2}}\,\frac{(r+s)^{m'}}{(r+s)^{2}}.
\end{align}
In two dimensions $H$ diverges logarithmically at large $|r|$: this is the collinear divergence
(large daughter separation $\leftrightarrow$ small relative momentum).

\textbf{Antisymmetric structure.} $H^{mm'}(s)=H^{m'm}(s)$ (substitute $r\to-r-s$), so the
contraction with $2(\delta_{mk'}\delta_{km'}-\delta_{km}\delta_{k'm'})$ vanishes identically,
at every $(\xi,y,y')$.

\textbf{Closed form of $H$ (the only regulated object).} In the scheme above,
\begin{align}\label{eq:H}
H^{mm'}(s)=-\frac{\pi}{\epsilon}\frac{\delta^{mm'}}{d}-\frac{\pi}{2}\,\delta^{mm'}\ln\frac{\mu^{2}s^{2}e^{2\gamma_E}}{4}-\pi\Big(\hat s^{m}\hat s^{m'}-\frac12\delta^{mm'}\Big)+O(\epsilon),
\end{align}
using $\int d^{2}\ell\,[e^{i\ell\cdot s}-\theta(\mu-|\ell|)]/\ell^{2}=-\pi\ln(\mu^{2}s^{2}e^{2\gamma_E}/4)$
and $\int d^{2}\ell\,e^{i\ell\cdot s}(\ell^{m}\ell^{m'}-\frac12\delta^{mm'}\ell^{2})/\ell^{4}=-\pi(\hat s^{m}\hat s^{m'}-\frac12\delta^{mm'})$
(equivalently ${\rm FT}[\delta-2\hat r\hat r]=\frac{4\pi}{q^{2}}(2\hat q\hat q-\delta)$).

\textbf{Contraction.} With $d_\perp=d$ the pole contracts to exactly $\Ph(\xi)$; the finite trace
gives $\Ph$ times the log; the traceless part survives only in the $d_\perp$ channel
($d_\perp\to2$ there, no pole). With $s^{2}=|y'-y|^{2}/\xi^{2}$:
\begin{align}\label{eq:I}
\boxed{\;I(\xi)=-\pi\,\Ph(\xi)\,G(\xi)\;+\;I_{\rm az}(\xi),\;}
\end{align}
\begin{align}
G(\xi)&\equiv\frac{1}{\xi^{2}}\int_{y,y'}e^{-ik\cdot(y'-y)/\xi}\Big[\frac{1}{\epsilon}+\ln\frac{\mu^{2}|y'-y|^{2}e^{2\gamma_E}}{4\xi^{2}}\Big]\Ac(y,y'),\label{eq:G}\\
I_{\rm az}(\xi)&\equiv\pi\,\frac{1-\xi}{\xi}\int_{y,y'}e^{-ik\cdot(y'-y)/\xi}\Big[\Ac(y,y')-2\Ac_{\hat s}(y,y')\Big].\label{eq:Iaz}
\end{align}
This is the note's Eq.~(50), derived here without leaving coordinate space.

\textbf{Remark (the $\xi^{2}$ inside the logarithm).} The collinear scale is $|w'-w|=|y'-y|/\xi$,
the separation of the observed gluons. Continuing $z$ at fixed parents instead of $r$ produces
$\ln\mu^{2}|y'-y|^{2}$ without the $\xi^{2}$, because $d^{d}z=\xi^{d}d^{d}r$. The two choices
differ by the finite term $+\pi\int d\xi\,\Ph(\xi)\ln\xi^{2}\,\Lc(k/\xi)/\xi^{2}$. Only the
$r$-continuation, which is the relative-momentum $\overline{\rm MS}$ scheme, matches the
$\overline{\rm MS}$ fragmentation function. Everything else agrees with the coordinate-space
results obtained earlier:
\begin{itemize}
\item the trace $-\pi[\frac1\epsilon+\ln]$ is your $\int_zK=-\pi L_\epsilon$ (with $2/\epsilon_{\rm yours}=1/\epsilon$);
\item $I_{\rm az}$ is exactly your bracket term $C\frac{1-\xi}{\xi}[\mathcal K\!\cdot\!\mathcal K'-2(\mathcal K\!\cdot\!\hat r)(\mathcal K'\!\cdot\!\hat r)]$.
\end{itemize}

\section{Longitudinal endpoints in coordinate space}
\textbf{$\xi\to1$ (target rapidity, JIMWLK).} At $\xi=1$, $y\to w$ and $y'\to w'$ and $s=y'-y$:
\begin{align}\label{eq:J}
\Jc(k)\equiv\lim_{\xi\to1}(1-\xi)I(\xi)=-\pi G(1)
=\int_{z,w,w'}e^{-ik\cdot(w'-w)}\frac{(w-z)\cdot(w'-z)}{(w-z)^{2}(w'-z)^{2}}\,\Ac(w,w')
=-\pi\!\int_{w,w'}\!e^{-ik\cdot(w'-w)}\Big[\frac1\epsilon+\ln\frac{\mu^{2}|w'-w|^{2}e^{2\gamma_E}}{4}\Big]\Ac(w,w').
\end{align}
\textbf{$\xi\to0$ (projectile rapidity, BFKL).} Since $\xi\Ph\to1$, the coefficient is
$-\pi G(0^{+})$. Put $y'=y+\xi t$, so $d^{2}y'=\xi^{2}d^{2}t$:
\begin{align}
G(\xi)=\int_{y}\int_{t}e^{-ik\cdot t}\Big[\frac1\epsilon+\ln\frac{\mu^{2}t^{2}e^{2\gamma_E}}{4}\Big]\Ac(y,y+\xi t)\xrightarrow{\xi\to0}\int_{y}\Ac(y,y)\int_{t}e^{-ik\cdot t}\ln t^{2}=-\frac{4\pi}{k^{2}}\int_{y}\Ac(y,y),
\end{align}
using $\int d^{2}t\,e^{-ik\cdot t}\ln t^{2}=-4\pi/k^{2}$ ($k\neq0$; the constant and the pole
multiply $\delta^{(2)}(k)$). Hence
\begin{align}\label{eq:B}
\Bc(k)\equiv\lim_{\xi\to0}\xi I(\xi)=-\pi G(0^{+})=\frac{4\pi^{2}}{k^{2}}\int_{z}\Ac(z,z)
=\int_{x,x',z,w,w'}e^{-ik\cdot(w'-w)}\frac{(w-z)\cdot(w'-z)}{(w-z)^{2}(w'-z)^{2}}\frac{(x-z)\cdot(x'-z)}{(x-z)^{2}(x'-z)^{2}}\big[U(z)-U(x')\big]^{ce'}\big[U(z)-U(x)\big]^{ce}\rho^{e'}\rho^{e}.
\end{align}
$\Bc$ has no collinear pole. $I_{\rm az}$ contributes to neither endpoint. Near $\xi\to0$ the
traceless bracket has vanishing angular average and its transform is
$\propto1/q^{2}$ at $q=k/\xi$, so $\xi I_{\rm az}=O(\xi^{2})$. Near $\xi\to1$ it carries $(1-\xi)$.

\textbf{Endpoint subtraction.} Use the exact decomposition
$\Ph(\xi)=\frac{1}{1-\xi}+\frac1\xi+\big(\xi(1-\xi)-2\big)$:
\begin{align}\label{eq:endpoints}
\int_{\xi_V}^{\xi_\Lambda}\!d\xi\,I(\xi)
=\Jc\ln\frac{1-\xi_V}{1-\xi_\Lambda}+\Bc\ln\frac{\xi_\Lambda}{\xi_V}
-\pi\!\int_{\xi_V}^{\xi_\Lambda}\!d\xi\Big[\frac{G(\xi)-G(1)}{1-\xi}+\frac{G(\xi)-G(0^{+})}{\xi}+\big(\xi(1-\xi)-2\big)G(\xi)\Big]
+\int_{\xi_V}^{\xi_\Lambda}\!d\xi\,I_{\rm az}(\xi),
\end{align}
with $\ln\frac{1-\xi_V}{1-\xi_\Lambda}=\ln\frac{k^{+}}{\Lambda}+\ln(1-\frac{k^{+}}{\vee})+O(\frac{\Lambda}{k^{+}})$
and $\ln\frac{\xi_\Lambda}{\xi_V}=\ln\frac{\vee}{k^{+}}+O(\frac{\Lambda}{k^{+}})$. The remaining
integrals converge at both ends, so their limits can be taken to $[0,1]$ at the cost of
$O(k^{+}/\vee,\Lambda/k^{+})$. (Your earlier finite-$\vee$ form, with subtractions
$\xi\Jc/(1-\xi)$ and $(1-\xi)\Bc/\xi$, is Eq.~\eqref{eq:endpoints} rearranged. It differs by
$-(1-\frac{k^{+}}{\vee})(\Jc+\Bc)$ and the $\ln(1-\frac{k^{+}}{\vee})\Jc$ term, which is exact.)

\section{Collinear pole, plus prescription, and the overlap}
Split $G=\frac{1}{\epsilon}h+G_{\rm fin}$ with
$h(\xi)\equiv\Lc(k/\xi)/\xi^{2}$, $h(0^{+})=0$ (since $\Lc(p)\sim p^{-4}$), $h(1)=\Lc(k)$. The
pole of Eq.~\eqref{eq:endpoints} is
\begin{align}
-\frac{\pi}{\epsilon}\Big\{\ln\frac{k^{+}}{\Lambda}\Lc(k)+\int_{0}^{1}\!d\xi\Big[\frac{h-h(1)}{1-\xi}+\frac{h}{\xi}+\big(\xi(1-\xi)-2\big)h\Big]\Big\}
=-\frac{\pi}{\epsilon}\Big\{\ln\frac{k^{+}}{\Lambda}\Lc(k)+\int_{0}^{1}\!\frac{d\xi}{\xi^{2}}\,\Ph_{+}(\xi)\,\Lc(k/\xi)\Big\},
\end{align}
where $\Ph_{+}=\frac{\xi}{(1-\xi)_{+}}+\frac{1-\xi}{\xi}+\xi(1-\xi)$ uses the \emph{standard}
plus distribution, acting on the full test function. The two integrands agree pointwise
(checked symbolically).
\begin{itemize}
\item The $\xi\to0$ end of the collinear convolution converges ($h\sim\xi^{2}$), so the
collinear and projectile-rapidity divergences do not overlap.
\item The double divergence $-\frac\pi\epsilon\ln\frac{k^{+}}{\Lambda}\Lc(k)$ is counted once,
inside $\Jc$. The plus prescription prevents counting it a second time.
\end{itemize}
In DGLAP language, with the note's normalisation ($\Nn_{\rm LO}=\alpha_s/4\pi^{4}$),
$-\frac{\pi\bar C}{\epsilon k^{+}}\int\frac{d\xi}{\xi^{2}}\Ph_{+}\Lc(k/\xi)
=-\frac{1}{\epsilon}\frac{\alpha_s}{2\pi}\int\frac{d\xi}{\xi^{3}}[2N_c\Ph_{+}(\xi)]\frac{d^{3}N_{\rm LO}}{d^{3}p}\big|_{p=k/\xi}$:
the real part of $P_{gg}\otimes{\rm LO}$, with $\xi^{-3}$ the fragmentation Jacobian.

\textbf{No UV divergence.} In coordinate space the short-distance points $z\to w$ and
$z\to w'$ ($r\to0$, $r\to-s$) are $1/|r|$ singularities, integrable in $d=2$. The only
transverse divergences of $H$ are at large $|r|$ (IR).

\section{The $\delta(1-\xi)$ completion and the coordinate-space running coupling}
$\Ph_{+}=\frac{P_{gg}}{2N_c}-\frac{11}{12}\delta(1-\xi)$ ($n_f=0$, $\beta_0=\frac{11N_c}{3}$;
momentum sum rule $\int_{0}^{1}\xi\Ph_{+}=-\frac{11}{12}$). Add and subtract:
\begin{align}
-\pi\!\int\!d\xi\,\Ph\,G
=-\pi\!\int\!d\xi\,\frac{P_{gg}(\xi)}{2N_c}\,G(\xi)\;+\;\pi\,\frac{11}{12}\,G(1)\;-\;\pi\ln\frac{k^{+}}{\Lambda}\,G(1).
\end{align}
\begin{itemize}
\item \emph{Fragmentation counterterm} at $\mu_F$ turns the first term into
$-\pi\int d\xi\frac{P_{gg}}{2N_c}\frac{1}{\xi^{2}}\int_{y,y'}e^{-ik\cdot(y'-y)/\xi}\ln\frac{\mu_F^{2}|y'-y|^{2}e^{2\gamma_E}}{4\xi^{2}}\Ac(y,y')$.
\item \emph{Coupling counterterm} at $\mu_R$ turns the second into
$+\pi\frac{11}{12}\int_{y,y'}e^{-ik\cdot(y'-y)}\ln\frac{\mu_R^{2}|y'-y|^{2}e^{2\gamma_E}}{4}\Ac(y,y')$,
i.e.\ $\frac{\alpha_s\beta_0}{4\pi}\ln\frac{\mu_R^{2}|y'-y|^{2}e^{2\gamma_E}}{4}$ inside the LO
coordinate integral. Combined with $\alpha_s(\mu_R)$ in the LO prefactor,
\begin{align}\label{eq:rc}
\frac{d^{3}N}{d^{3}k}\supset\frac{\Nn'}{k^{+}}\int_{y,y'}\alpha_s\Big(\frac{4e^{-2\gamma_E}}{|y'-y|^{2}}\Big)\,e^{-ik\cdot(y'-y)}\,\Ac(y,y'),\qquad\Nn'=\frac{1}{4\pi^{4}}:
\end{align}
the coupling runs at the parent separation $|y'-y|$ \emph{inside} the Fourier integral, not at
$k^{2}$.
\item The third term, $-\pi\ln\frac{k^{+}}{\Lambda}G(1)=\ln\frac{k^{+}}{\Lambda}\Jc$, carries the
pole $-\frac\pi\epsilon\ln\frac{k^{+}}{\Lambda}\Lc(k)$. It must cancel against the $z\to w$ UV poles
of the soft-gluon-scattering ($J_LU(z)J_R$-type) terms, which are not part of this row.
\end{itemize}
This term has no UV pole of its own. The add-and-subtract is the required bookkeeping once
scaleless virtual integrals are set to zero: the final-state self-energy has the would-be value
$\pi[\ln\frac{k^{+}}{\Lambda}-\frac{11}{12}]\Lc(k)(\frac{1}{\epsilon_{\rm IR}}-\frac{1}{\epsilon_{\rm UV}})\equiv0$.

\section{Net-colour large-distance IR divergence (coordinate space)}
With $U\to1$ outside the target, at large $|y|$
\begin{align}
A^{c}_{i}(y)\xrightarrow{|y|\to\infty}-\frac{y_i}{y^{2}}\,Q^{c},\qquad
Q^{c}=\int_x\big[\delta^{ce}-U^{ce}(x)\big]\rho^{e}(x),
\end{align}
generically nonzero, even for a colour-neutral projectile. Then
$\int_{|y|<R}d^{2}y\,\Ac(y,y)=2\pi Q^{c}Q^{c}\ln R+{\rm const}$, so $\Bc$ diverges
logarithmically at large $|z|$. In the same scheme, with
$\Lc_C(k)\equiv4\pi^{2}Q^{c}Q^{c}/k^{2}$ (the Coulomb tail of $\Lc$),
\begin{align}
\Bc=\hat\Bc-\frac{\pi}{\epsilon}\Lc_C(k),\qquad
\Jc=\hat\Jc_{\rm fin}-\frac{\pi}{\epsilon}\Lc(k)-\frac{\pi}{\epsilon}\Lc_C(k),\qquad
I=\hat I-\frac{\pi}{\epsilon}\frac{\Ph}{\xi^{2}}\Lc(k/\xi)-\frac{\pi}{\epsilon}\Ph\,\Lc_C(k),
\end{align}
(the large-$|y|,|y'|$ region of the coordinate integrals), and, using
$\int_{0}^{1}[\Ph-\frac{1}{1-\xi}-\frac1\xi]=\int_{0}^{1}(\xi(1-\xi)-2)=-\frac{11}{6}$,
\begin{align}
\frac{d^{3}N}{d^{3}k}\supset-\frac{\pi\bar C}{\epsilon\,k^{+}}\,\Lc_C(k)\Big[\ln\frac{k^{+}}{\Lambda}+\ln\frac{\vee}{k^{+}}-\frac{11}{6}\Big].
\end{align}
With a large-distance cutoff, $-1/\epsilon\leftrightarrow\ln R^{2}$ up to a constant. This pole
cannot cancel within this row.

\section{Result}
\textbf{Compact form} (exact up to $O(\Lambda/k^{+},k^{+}/\vee)$, all transverse integrals in
coordinate space):
\begin{align}\label{eq:result}
\boxed{\;\frac{d^{3}N}{d^{3}k}=\Nn\Big\{\ln\frac{k^{+}}{\Lambda}\,\Jc(k)+\ln\frac{\vee}{k^{+}}\,\Bc(k)
-\pi\!\int_{0}^{1}\!d\xi\Big[\frac{G(\xi)-G(1)}{1-\xi}+\frac{G(\xi)-G(0^{+})}{\xi}+\big(\xi(1-\xi)-2\big)G(\xi)\Big]
+\int_{0}^{1}\!d\xi\,I_{\rm az}(\xi)\Big\}\;}
\end{align}
with $\Nn=\frac{1}{(2\pi)^{3}}\frac{g^{4}N_c}{8\pi^{5}}\frac{1}{k^{+}}$, and $G$, $I_{\rm az}$, $\Jc=-\pi G(1)$, $\Bc=-\pi G(0^{+})$ from
Eqs.~\eqref{eq:G}, \eqref{eq:Iaz}, \eqref{eq:J}, \eqref{eq:B}.

\textbf{All poles explicit} (coordinate version of the note's Eq.~(65)):
\begin{align}
\frac{d^{3}N}{d^{3}k}=\Nn\Big\{&\underbrace{-\frac{\pi}{\epsilon}\int_{0}^{1}\frac{d\xi}{\xi^{2}}\Big[\Ph_{+}+\frac{11}{12}\delta(1-\xi)\Big]\Lc(k/\xi)}_{\text{full }P_{gg}:\ \overline{\rm MS}\text{ fragmentation c.t.}}
\underbrace{+\frac{\pi}{\epsilon}\frac{11}{12}\Lc(k)}_{\text{cancels }\alpha_s\text{ c.t.}}
\underbrace{-\frac{\pi}{\epsilon}\ln\frac{k^{+}}{\Lambda}\Lc(k)}_{\text{UV of target evolution}}
\underbrace{-\frac{\pi}{\epsilon}\Lc_C(k)\Big[\ln\frac{k^{+}}{\Lambda}+\ln\frac{\vee}{k^{+}}-\frac{11}{6}\Big]}_{\text{net-colour IR}}\notag\\
&+\ln\frac{k^{+}}{\Lambda}\,\hat\Jc_{\rm fin}(k)+\ln\frac{\vee}{k^{+}}\,\hat\Bc(k)+\hat F(k;\mu)\Big\},
\end{align}
with the coordinate-space finite pieces
\begin{align}
\Jc_{\rm fin}(k)&=-\pi\int_{w,w'}e^{-ik\cdot(w'-w)}\ln\frac{\mu^{2}|w'-w|^{2}e^{2\gamma_E}}{4}\,\Ac(w,w'),\qquad
\Bc(k)=\frac{4\pi^{2}}{k^{2}}\int_{z}\Ac(z,z),\\
F(k;\mu)&=\int_{0}^{1}d\xi\Big[I_{\rm fin}(\xi)-\frac{\Jc_{\rm fin}(k)}{1-\xi}-\frac{\Bc(k)}{\xi}\Big],\\
I_{\rm fin}(\xi)&=-\frac{\pi}{\xi^{2}}\int_{y,y'}e^{-ik\cdot(y'-y)/\xi}\Big\{\Ph(\xi)\ln\frac{\mu^{2}|y'-y|^{2}e^{2\gamma_E}}{4\xi^{2}}\,\Ac(y,y')+2\xi(1-\xi)\Big[\Ac_{\hat s}(y,y')-\frac12\Ac(y,y')\Big]\Big\},
\end{align}
and hats denoting removal of the net-colour (Coulomb-tail) pieces. After the $\overline{\rm MS}$
counterterms, the $\beta_0$ share becomes the coordinate-space running coupling of
Eq.~\eqref{eq:rc}.

\section{Checks}
\begin{itemize}
\item Change of variables: Jacobian $1/\xi^{2}$ and $w'-w=(y'-y)/\xi$ (symbolic);
$H^{mm'}=H^{m'm}$ kills the antisymmetric structure exactly.
\item $\Ph-\frac1{1-\xi}-\frac1\xi=\xi(1-\xi)-2$, $\int_{0}^{1}(\xi(1-\xi)-2)=-\frac{11}{6}$, and
the subtracted collinear integrand equals the standard-plus integrand pointwise (symbolic);
$\int_{0}^{1}\xi\Ph_{+}=-\frac{11}{12}$.
\item $-\pi G(1)=\Jc$ is the dimensionally regularised $\int_zK_{ww'z}$; $-\pi G(0^{+})=\Bc$
follows from ${\rm FT}[\ln t^{2}]=-4\pi/k^{2}$. These are the JIMWLK and BFKL coefficients
obtained earlier, with the same signs and prefactors.
\item $I_{\rm az}$ reproduces the earlier finite bracket term. Its momentum-space image
$(4\pi/q^{2})(2\hat q\hat q-\delta)$ was verified distributionally to 8 digits.
\item $I(\xi)$ of Eq.~\eqref{eq:I} is the note's Eq.~(50). The note checked it against the
momentum-space disc-subtracted integrand to $9\times10^{-7}$ and $3\times10^{-8}$.
\end{itemize}
\end{document}
